% Generated draft: requires mathematical and source review.
\documentclass[11pt]{article}
\usepackage[margin=25mm]{geometry}
\usepackage{fontspec}
\setmainfont{Latin Modern Roman}
\usepackage{amsmath,amssymb,mathtools,unicode-math}
\setmathfont{Latin Modern Math}
\usepackage{enumitem,needspace,xurl,hyperref}
\setlength{\parindent}{0pt}
\setlength{\parskip}{5pt}
\setlength{\emergencystretch}{3em}
\newcommand{\N}{\mathbb{N}}
\newcommand{\Z}{\mathbb{Z}}
\newcommand{\Q}{\mathbb{Q}}
\newcommand{\R}{\mathbb{R}}
\newcommand{\C}{\mathbb{C}}
\newcommand{\F}{\mathbb{F}}
\newcommand{\A}{\mathbb{A}}
\newcommand{\PP}{\mathbb{P}}
\newcommand{\E}{\mathbb{E}}
\newcommand{\eps}{\varepsilon}
\newcommand{\dd}{\,\mathrm{d}}
\begin{document}
\section*{UnsolvedMath ID: 20002912}
\textbf{Source problem code:} catalogue entry
\subsection{Definitions and mathematical statement}
The AIM workshop question asks for a good definition of level-restricted $q,t$-Kostka numbers, asks whether such a definition makes sense for root systems other than type $A$, and asks for its relationship with $k$-$q,t$-Kostka numbers. Here $q,t$ denote deformation parameters. A level restriction is additional representation-theoretic data; the source does not specify its normalization, the indexing set, or axioms characterizing a ``good'' definition. The parameter called $k$ in the last question must not be identified with the level without an explicit convention.

The supplied research aid concerns the generalized version with a sequence $R=(R_1,\ldots,R_m)$ of rectangular partitions. Write $s_\nu$ for a Schur function and $[s_\lambda]f$ for its coefficient in $f$. At stable level, the ungraded tensor-product multiplicity associated with this sequence is
\[
 c_{\lambda,R}=[s_\lambda]\prod_{j=1}^m s_{R_j}.
\]
The aid gives the following obstruction to forgetting the block structure. In type $A$, with at least two available rows, take
\[
 R_{\mathrm{col}}=((1,1)),\qquad
 R_{\mathrm{split}}=((1),(1)),\qquad \lambda=(2).
\]
The first sequence consists of one two-box column, whereas the second consists of two one-box rectangles. Forgetting the rectangle separators and retaining only their row lengths gives $\mu=(1,1)$ in both cases. Nevertheless,
\[
 [s_{(2)}]s_{(1,1)}=0,\qquad
 s_{(1)}^2=s_{(2)}+s_{(1,1)},\qquad
 [s_{(2)}]s_{(1)}^2=1.
\]
Consequently, no family depending only on $(\lambda,\mu,\ell)$, where $\ell$ is the level and $\mu$ is precisely this separator-forgetting encoding, can specialize to all generalized level-restricted multiplicities: choose a common stable level for these two sequences, and its value at $q=t=1$ would have to be both $0$ and $1$. A proposed two-parameter refinement must retain $R$, or restrict to a specified canonical choice $R=R(\mu)$. This does not obstruct ordinary Kostka families with a fixed interpretation of their content index.

This obstruction leaves the original construction problem open in the supplied record, including its non-type-$A$ part. The aid proposes comparing modern positive $k$-$q,t$ coefficients with crystal or fusion coefficients only after the rectangle subfamily, index correspondence, and normalization are specified. It singles out an explicitly normalized $Q=0$ boundary identity as the remaining comparison test. No formula for that identity, definition of $Q$, or relation between $Q$ and the original parameters $q,t$ is supplied, so an exact boundary equality cannot be asserted here.
\subsection{Short English statement}
Define a level-restricted two-parameter Kostka theory, determine its scope beyond type $A$, and compare it with $k$-$q,t$-Kostka numbers. The supplied partial result shows that a universal generalized theory cannot discard how its rectangles are grouped: one column and two separate boxes have different stable multiplicities despite identical ungrouped row data.
\subsection{Sources}
\begin{itemize}
\item Local UnsolvedMath registry, permanent ID $20002912$, source code catalogue entry; canonical location \texttt{aim-special-functions-notes.json}, \texttt{notes[16]}.
\item American Institute of Mathematics, workshop \emph{Generalized Kostka polynomials}, problem (17), \url{https://aimath.org/WWN/kostka/prob2.pdf}. The supplied extraction contains no proposer, remarks, or literature entries.
\item Research attempt 1 and research summary attached to the local record: rectangle-data obstruction and proposed normalized $Q=0$ comparison test. The record labels this partial progress and labels the attempt proved with low novelty confidence, while explicitly warning that it is a machine-generated, non-peer-reviewed research aid requiring independent verification before citation.
\end{itemize}
\end{document}
